% ==============================================================================
% 09_threat_intelligence.tex — UE SIS589
% ==============================================================================
\chapter{Cyber Threat Intelligence}
\label{chap:cti}

\epigraph{\itshape « Knowing your enemy is the first condition of victory. »}
         {--- Sun Tzu, \textit{L'Art de la Guerre}}

\begin{boxobjectifs}
\begin{enumerate}
  \item Définir la CTI et ses quatre niveaux~: stratégique, opérationnel,
        tactique, technique.
  \item Appliquer les modèles structurants~: Cyber Kill Chain, Diamond Model,
        MITRE ATT\&CK.
  \item Distinguer et exploiter les IoC, IoA et TTP.
  \item Décrire les formats d'échange STIX/TAXII et les plateformes
        MISP/OpenCTI.
  \item Intégrer la CTI dans les règles SIEM et les procédures de réponse.
  \item Évaluer la crédibilité d'une source de threat intelligence.
\end{enumerate}
\end{boxobjectifs}

% ==============================================================================
\section{Définition et niveaux de la CTI}
\label{sec:cti-def}
% ==============================================================================

\begin{boxdef}
\textbf{Cyber Threat Intelligence (CTI)~:} Renseignement produit à partir
de la collecte, du traitement et de l'analyse d'informations sur les menaces
cybernétiques. La CTI est \textit{actionnelle}~: elle répond à une question
décisionnelle spécifique et est fournie au bon acteur, au bon moment, sous
le bon format. Un flux d'IoC brut n'est pas de la CTI~; c'est une donnée
brute. La CTI est le produit de son analyse.
\end{boxdef}

La CTI opère à quatre niveaux distincts, correspondant à des audiences
et des horizons temporels différents~:

\begin{table}[H]
\centering\small
\caption{Niveaux de la Cyber Threat Intelligence}
\label{tab:cti-niveaux}
\rowcolors{2}{TableauPair}{TableauImpair}
\begin{tabular}{p{2.8cm}p{3cm}p{3.5cm}p{5cm}}
\toprule
\tableauHeader{Niveau} & \tableauHeader{Audience} &
\tableauHeader{Horizon} & \tableauHeader{Contenu type} \\
\midrule
\textbf{Stratégique} &
Direction, COMEX &
Long terme &
Tendances géopolitiques, secteurs ciblés, impacts business \\
\textbf{Opérationnel} &
RSSI, responsables SOC &
Moyen terme &
Campagnes actives, groupes d'attaque, intentions \\
\textbf{Tactique} &
Analystes N2/N3 &
Court terme &
TTPs des adversaires, techniques d'attaque, playbooks adversaires \\
\textbf{Technique} &
Analystes N1, SIEM &
Immédiat &
IoC~: hashes, IPs, domaines, URLs, règles YARA/Sigma \\
\bottomrule
\end{tabular}
\end{table}

\begin{boiteRemarque}{CTI versus données brutes}
Un flux d'adresses IP malveillantes est une \textbf{donnée brute}.
La même liste, filtrée pour ne conserver que les IPs qui communiquent
avec des secteurs similaires au vôtre, datée de moins de 30~jours,
avec un niveau de confiance documenté et un contexte de campagne~:
c'est de la \textbf{CTI technique actionnelle}. La valeur ajoutée de
l'analyste CTI réside dans ce traitement.
\end{boiteRemarque}

% ==============================================================================
\section{Indicateurs~: IoC, IoA, TTP}
\label{sec:indicateurs}
% ==============================================================================

\begin{boxdef}
\textbf{IoC (\textit{Indicator of Compromise})~:} Artefact observable dans
un système ou un réseau, indiquant une compromission passée ou en cours.
Un IoC est \textbf{réactif} (postérieur à l'attaque) et a une durée de vie
limitée (les attaquants changent leurs IPs, domaines et hashes régulièrement).
Exemples~: hash SHA-256 d'un malware, IP d'un serveur C2, domaine DGA,
chemin de fichier suspect.
\end{boxdef}

\begin{boxdef}
\textbf{IoA (\textit{Indicator of Attack})~:} Indicateur de comportement
suspect \textbf{en cours d'attaque}, indépendamment des artefacts.
Un IoA est \textbf{proactif}~: il détecte l'intention ou l'action avant
la compromission complète. Plus durable qu'un IoC (les comportements
changent moins vite que les artefacts).
Exemples~: exécution de PowerShell avec paramètre encodé, création de
compte admin hors des heures ouvrables, connexion RDP depuis un pays
inhabituel.
\end{boxdef}

\begin{boxdef}
\textbf{TTP (\textit{Tactics, Techniques and Procedures})~:}
Comportements, méthodes et procédures caractéristiques d'un acteur
de la menace. Niveau le plus durable et le plus difficile à changer
pour un attaquant. Les TTPs définissent le \og{}style opératoire\fg{}
d'un groupe (APT29, Lazarus, etc.). Référence normative~: MITRE ATT\&CK.
\end{boxdef}

\begin{boiteRemarque}{La pyramide de la douleur}
David Bianco (\textit{Pyramid of Pain}, 2013) hiérarchise les IoC
par la difficulté qu'ils imposent à l'attaquant lorsqu'ils sont
bloqués~: hashes (trivial à changer)~$<$~IPs~$<$~domaines~$<$~artefacts
réseau~$<$~outils~$<$~TTPs (très difficile à changer). Bloquer des TTPs
est la stratégie défensive la plus durable.
\end{boiteRemarque}

% ==============================================================================
\section{Modèles de structuration}
\label{sec:cti-modeles}
% ==============================================================================

\subsection{Cyber Kill Chain (Lockheed Martin, 2011)}

\begin{boxdef}
\textbf{Cyber Kill Chain~:} Modèle décrivant les sept phases séquentielles
d'une attaque avancée. Rompre la chaîne à n'importe quel maillon empêche
la réalisation de l'objectif final.
\end{boxdef}

\begin{enumerate}
  \item \textbf{Reconnaissance~:} collecte d'informations sur la cible (OSINT,
        scan, phishing ciblé). Correspondance ATT\&CK~: TA0043.
  \item \textbf{Weaponization~:} création du vecteur d'attaque (exploit +
        payload). Aucune interaction cible.
  \item \textbf{Delivery~:} livraison du vecteur (email de phishing, clé USB,
        exploitation d'un service exposé). ATT\&CK~: TA0001.
  \item \textbf{Exploitation~:} déclenchement de la vulnérabilité sur le
        système cible. ATT\&CK~: TA0002.
  \item \textbf{Installation~:} dépôt d'un accès persistant (implant,
        backdoor, RAT). ATT\&CK~: TA0003.
  \item \textbf{Command \& Control (C2)~:} établissement d'un canal de
        communication entre l'implant et l'infrastructure de l'attaquant.
        ATT\&CK~: TA0011.
  \item \textbf{Actions on Objectives~:} exfiltration, chiffrement
        (ransomware), sabotage. ATT\&CK~: TA0009, TA0040.
\end{enumerate}

La Kill Chain est utile pour les discussions stratégiques mais souffre
d'une limite~: elle est séquentielle et ne modélise pas les attaquants
qui sont déjà à l'intérieur ou qui opèrent en cycles non linéaires.

\subsection{Diamond Model}

\begin{boxdef}
\textbf{Diamond Model (Caltagirone, Pendergast, Betz, 2013)~:}
Modèle d'analyse des incidents structurant chaque événement malveillant
autour de quatre sommets~: \textbf{Adversaire} (qui), \textbf{Infrastructure}
(comment~: IPs, domaines), \textbf{Capacité} (quoi~: outils, exploits)
et \textbf{Victime} (cible). Les méta-fonctionnalités (horodatage,
phase de la Kill Chain, orientation socio-politique) complètent le modèle.
\end{boxdef}

L'intérêt du Diamond Model est la \textbf{pivotisation}~: si l'on
connaît un sommet (ex. une IP d'infrastructure), on peut inférer les
autres sommets~; si l'on connaît une capacité (ex. un outil), on peut
retrouver d'autres victimes ou l'adversaire associé.

\subsection{MITRE ATT\&CK}

\begin{boxdef}
\textbf{MITRE ATT\&CK (\textit{Adversarial Tactics, Techniques and Common
Knowledge})~:} Base de connaissance mondiale des TTPs d'attaquants réels,
structurée en matrices par environnement (Enterprise, Mobile, ICS).
14~tactiques (colonnes), $>$200~techniques, $>$400~sous-techniques.
Mise à jour régulière. Standard de facto pour l'annotation des incidents,
l'ingénierie de détection et la mesure de couverture défensive.
\end{boxdef}

Les 14~tactiques de la matrice Enterprise~:
TA0043~Reconnaissance, TA0042~Resource Development,
TA0001~Initial Access, TA0002~Execution, TA0003~Persistence,
TA0004~Privilege Escalation, TA0005~Defense Evasion,
TA0006~Credential Access, TA0007~Discovery, TA0008~Lateral Movement,
TA0009~Collection, TA0011~C\&C, TA0010~Exfiltration,
TA0040~Impact.

% ==============================================================================
\section{Formats d'échange~: STIX et TAXII}
\label{sec:stix-taxii}
% ==============================================================================

\begin{boxdef}
\textbf{STIX (\textit{Structured Threat Information eXpression})~:}
Format JSON standardisé (OASIS, v2.1) pour représenter des objets
de threat intelligence~: indicateurs (\textit{Indicator}),
acteurs (\textit{Threat-Actor}), campagnes (\textit{Campaign}),
TTPs (\textit{Attack-Pattern}), outils (\textit{Tool}),
vulnérabilités (\textit{Vulnerability}), et leurs relations.
\end{boxdef}

\begin{boxdef}
\textbf{TAXII (\textit{Trusted Automated eXchange of Intelligence
Information})~:} Protocole HTTP(S) standardisé (OASIS) pour la
transmission automatique d'objets STIX entre plateformes de threat
intelligence. Deux services~: \textit{Collection} (dépôt d'objets)
et \textit{Channel} (publication push).
\end{boxdef}

\begin{lstlisting}[style=StyleCode, language=YAML,
  caption={Objet STIX~2.1 --- indicateur IP malveillante},
  label={lst:stix-indicator}]
{
  "type"             : "indicator",
  "spec_version"     : "2.1",
  "id"               : "indicator--a6e08ad4-9e2f-4c0d-b4f2-3e7c1d8b2a3e",
  "created"          : "2026-04-12T05:00:00Z",
  "modified"         : "2026-04-12T05:00:00Z",
  "name"             : "IP C2 incident LiZbiZ-SIS 12/04/2026",
  "description"      : "Adresse IP utilisée comme serveur C2 lors de l'incident",
  "pattern"          : "[ipv4-addr:value = '185.220.101.5']",
  "pattern_type"     : "stix",
  "valid_from"       : "2026-04-12T03:17:00Z",
  "indicator_types"  : ["malicious-activity", "compromised"],
  "confidence"       : 95,
  "kill_chain_phases": [{
    "kill_chain_name" : "mitre-attack",
    "phase_name"      : "command-and-control"
  }],
  "labels"           : ["incident:INC-2026-0412-001", "tlp:amber"]
}
\end{lstlisting}

% ==============================================================================
\section{Plateformes de CTI~: MISP et OpenCTI}
\label{sec:cti-plateformes}
% ==============================================================================

\subsection{MISP (\textit{Malware Information Sharing Platform})}

\begin{boxdef}
\textbf{MISP~:} Plateforme open-source de partage d'informations sur
les menaces. Supporte les formats STIX, OpenIOC, CSV. Permet la création
d'événements (incidents), le tagging (TLP, PAP, MITRE ATT\&CK),
la corrélation automatique entre événements, et la fédération avec
d'autres instances MISP (partage sectoriel, national, CIRCL).
\end{boxdef}

Fonctions clés pour le SOC~:
\begin{itemize}
  \item Alimenter le SIEM en IoC (export CSV/JSON vers Logstash).
  \item Corréler un hash ou une IP observée lors d'un incident avec
        les événements du passé.
  \item Partager les IoC de l'incident avec les partenaires (secteur,
        CERT-MU Cameroun).
  \item Consommer les flux CTI externes (Emerging Threats, AlienVault OTX,
        abuse.ch, VirusTotal).
\end{itemize}

\begin{lstlisting}[style=StyleTerminal,
  caption={MISP API --- recherche et export d'IoC},
  label={lst:misp-api}]
# ─── Rechercher une IP dans MISP (PyMISP) ────────────────────────────────────
python3 << 'EOF'
from pymisp import PyMISP

misp = PyMISP('https://misp.lizbiz.local', 'API_KEY', False)

# Rechercher si l'IP C2 est connue dans MISP
result = misp.search(value='185.220.101.5', type_attribute='ip-dst')
for event in result:
    print(f"Événement : {event['Event']['info']}")
    print(f"Date      : {event['Event']['date']}")
    print(f"Tags      : {[t['name'] for t in event['Event'].get('Tag', [])]}")
EOF

# ─── Exporter les IoC en CSV pour le SIEM (derniers 7 jours) ─────────────────
curl -s -H "Authorization: API_KEY" \
  "https://misp.lizbiz.local/attributes/restSearch/returnFormat:csv/\
type:ip-dst/last:7d/to_ids:1" \
  -o /etc/logstash/ioc_ips.csv

# ─── Créer un événement MISP depuis l'incident ────────────────────────────────
python3 << 'EOF'
from pymisp import PyMISP, MISPEvent, MISPAttribute

misp = PyMISP('https://misp.lizbiz.local', 'API_KEY', False)
event = MISPEvent()
event.info = "Incident LiZbiZ-SIS 12/04/2026 - Brute force SSH + C2"
event.threat_level_id = 1  # High
event.distribution = 1     # This community only

# Ajouter les IoC
for ip in ['203.0.113.42', '185.220.101.5']:
    event.add_attribute('ip-dst', ip, to_ids=True)
event.add_attribute('sha256',
  '3a7d4f8bc2e910...', comment='lp.sh payload', to_ids=True)

misp.add_event(event)
print("Événement créé avec succès")
EOF
\end{lstlisting}

\subsection{OpenCTI}

\begin{boxdef}
\textbf{OpenCTI~:} Plateforme open-source de gestion et de partage de
la CTI (Filigran, anciennement Luatix). Nativement orienté STIX~2.1.
Interface graphique riche pour la visualisation des relations entre
acteurs, campagnes, TTPs et indicateurs. Connecteurs vers MISP,
VirusTotal, AlienVault OTX, Shodan, Mandiant, CIRCL.
\end{boxdef}

Complémentarité MISP / OpenCTI~:
\begin{itemize}
  \item \textbf{MISP~:} optimisé pour la gestion opérationnelle des IoC
        et le partage communautaire.
  \item \textbf{OpenCTI~:} optimisé pour l'analyse stratégique/opérationnelle,
        la modélisation des acteurs et la visualisation des graphes de
        relations.
  \item Les deux s'interfacent via STIX~2.1 / API.
\end{itemize}

% ==============================================================================
\section{Intégration CTI dans le SOC}
\label{sec:cti-integration}
% ==============================================================================

\subsection{Enrichissement automatique du SIEM}

Le cycle d'intégration CTI $\to$ SIEM suit quatre étapes~:

\begin{enumerate}
  \item \textbf{Collecte CTI.} Abonnement aux flux (MISP, OTX, abuse.ch,
        Emerging Threats). Export périodique des IoC actifs (IP, domaines,
        hashes) en CSV/JSON.
  \item \textbf{Injection dans le pipeline.} Logstash translate plugin
        ou Wazuh CDB (Custom Database) pour le lookup en temps réel lors
        du parsing de chaque événement.
  \item \textbf{Alerte contextualisée.} Si un événement matche un IoC CTI,
        une alerte enrichie est générée avec le contexte~:
        nom de la campagne, groupe attributé, niveau de confiance.
  \item \textbf{Feedback.} Les IoC observés lors d'un incident sont
        réintégrés dans la plateforme CTI (boucle fermée).
\end{enumerate}

\begin{lstlisting}[style=StyleTerminal,
  caption={Intégration MISP $\to$ Wazuh CDB (Custom DB)},
  label={lst:misp-wazuh}]
# ─── Exporter les IoC de MISP vers une base CDB Wazuh ────────────────────────
# Format CDB Wazuh : une entrée par ligne "valeur:"
curl -s -H "Authorization: $MISP_KEY" \
  "https://misp.lizbiz.local/attributes/restSearch/\
returnFormat:csv/type:ip-dst/last:24h/to_ids:1" | \
  awk -F',' 'NR>1{print $2":"}' > \
  /var/ossec/etc/lists/malicious_ips.cdb

# Compiler la base CDB
/var/ossec/bin/ossec-makelists

# ─── Règle Wazuh : alerte si IP source dans la liste CTI ─────────────────────
# Dans /var/ossec/etc/rules/local_rules.xml :
# <rule id="100300" level="12">
#   <if_group>syslog,authentication</if_group>
#   <list field="srcip" lookup="match_key">etc/lists/malicious_ips</list>
#   <description>CTI Match: IP source dans base MISP (IOC connu)</description>
#   <mitre><id>T1071</id></mitre>
#   <group>cti_match,requires_triage</group>
# </rule>
\end{lstlisting}

\subsection{Threat Intel Feeds recommandés}

\begin{table}[H]
\centering\small
\caption{Sources de CTI technique gratuites et payantes}
\label{tab:cti-feeds}
\rowcolors{2}{TableauPair}{TableauImpair}
\begin{tabular}{p{3.5cm}p{3cm}p{3cm}p{4cm}}
\toprule
\tableauHeader{Source} & \tableauHeader{Type} &
\tableauHeader{Coût} & \tableauHeader{Contenu} \\
\midrule
AlienVault OTX      & STIX, CSV      & Gratuit   & IPs, domaines, hashes, URLs \\
abuse.ch (Feodo, MalwareBazaar) & CSV & Gratuit & C2 botnets, malwares \\
Emerging Threats    & Snort/Suricata & Gratuit + Pro & Règles IDS \\
CIRCL (MISP feeds)  & MISP           & Gratuit   & IoC gouvernement LU \\
VirusTotal          & API            & Freemium  & Hashes, URLs, domaines \\
Shodan              & API            & Freemium  & Empreintes infrastructure \\
Mandiant / Recorded Future & STIX  & Payant    & CTI opérationnelle avancée \\
\bottomrule
\end{tabular}
\end{table}

% ==============================================================================
\section{Évaluation et crédibilité des sources}
\label{sec:cti-credibilite}
% ==============================================================================

Tous les IoC ne se valent pas. Un IoC doit être évalué sur cinq
dimensions avant d'être intégré dans la production~:

\begin{enumerate}
  \item \textbf{Source.} Qui a produit l'IoC~? (CERT gouvernemental,
        éditeur sécurité, communauté anonyme, flux automatisé).
  \item \textbf{Confiance (\textit{confidence}).} Niveau de certitude
        que l'IoC est valide (0--100~\%). Un IoC à faible confiance
        ne doit pas déclencher un blocage automatique.
  \item \textbf{Fraîcheur (\textit{freshness}).} Un IoC vieux de 6~mois
        est probablement obsolète (infrastructure attaquant changée).
        Durée de vie recommandée~: 30~jours pour les IPs,
        90~jours pour les domaines, indéterminé pour les hashes.
  \item \textbf{Contexte.} L'IoC est-il pertinent pour votre secteur,
        votre géographie, votre architecture~?
  \item \textbf{Faux positifs.} Un IoC partagé trop largement peut
        inclure des IPs légitimes (CDN, cloud providers). Valider
        systématiquement avant blocage.
\end{enumerate}

\begin{boxdef}
\textbf{TLP (\textit{Traffic Light Protocol})~:} Système de classification
du partage de l'information en CTI~:
\textbf{TLP:RED}~=~destinataire uniquement~;
\textbf{TLP:AMBER}~=~organisation et partenaires de confiance~;
\textbf{TLP:GREEN}~=~communauté sectorielle~;
\textbf{TLP:CLEAR}~=~partage libre.
\end{boxdef}

% ==============================================================================
\section{CTI dans la réponse aux incidents}
\label{sec:cti-ir}
% ==============================================================================

La CTI enrichit chaque phase de la réponse~:

\begin{description}
  \item[\textbf{Préparation~:}] Modéliser les adversaires probables
        (secteur e-commerce, présence africaine). Configurer les
        playbooks en tenant compte des TTPs documentés.
  \item[\textbf{Détection~:}] Enrichissement automatique des alertes SIEM
        avec le contexte CTI. L'alerte \og{}brute force SSH depuis
        203.0.113.42\fg{} devient \og{}brute force SSH depuis IP appartenant
        au groupe FIN7 (Carbanak), campagne active mars--avril~2026\fg{}.
  \item[\textbf{Investigation~:}] Le Diamond Model guide la pivotisation~:
        de l'IP C2 vers les autres victimes, vers l'infrastructure,
        vers les outils, vers le groupe.
  \item[\textbf{Éradication~:}] Les TTPs documentés permettent de savoir
        quels autres vecteurs de persistance chercher, au-delà
        de ceux déjà découverts.
  \item[\textbf{RETEX~:}] Les IoC et TTPs de l'incident enrichissent
        la plateforme CTI, bénéficiant aux organisations partenaires.
\end{description}

\begin{boxlizbiz}
\textbf{LiZbiZ-SIS~: application CTI à l'incident du 12/04/2026}

Après l'incident, l'analyste N3 a pivotisé depuis l'IP C2
185.220.101.5 dans MISP~: cette IP était déjà référencée dans
3~événements liés au groupe \textit{ShadowFin} (acteur spécialisé
dans les PME e-commerce africaines). Les outils utilisés (wget + bash
+ C2 HTTPS) correspondent aux TTPs documentés de ce groupe~:
T1059.004, T1071.001, T1105. Cette information a permis de diriger
l'investigation vers d'autres artefacts typiques de ce groupe
(fichiers .lnk dans \%TEMP\%, clés de registre \texttt{WindowsUpdate})
qui n'auraient pas été cherchés sans le contexte CTI.
\end{boxlizbiz}

% ==============================================================================
\section*{Résumé}
% ==============================================================================

\begin{boxresume}
\begin{itemize}
  \item CTI~: 4~niveaux (stratégique, opérationnel, tactique, technique).
        La CTI est actionnelle~; un flux d'IoC brut n'est pas de la CTI.
  \item IoC (réactif, court terme)~$<$~IoA (comportemental)~$<$~TTP
        (durable). Bloquer les TTPs~: la stratégie la plus efficace
        (pyramide de la douleur).
  \item Modèles~: Kill Chain (7~phases séquentielles), Diamond Model
        (Adversaire / Infrastructure / Capacité / Victime), MITRE ATT\&CK
        (14~tactiques, $>$200~techniques).
  \item Formats~: STIX~2.1 (représentation), TAXII (transport).
  \item Plateformes~: MISP (partage opérationnel des IoC),
        OpenCTI (analyse stratégique et graphe de relations).
  \item Intégration SOC~: enrichissement SIEM via CDB Wazuh ou Logstash
        translate~; boucle CTI fermée.
  \item Crédibilité IoC~: source, confiance, fraîcheur, contexte, FP.
\end{itemize}
\end{boxresume}

\begin{boxexo}
\textbf{Ex.~9.1 — Classification des indicateurs.}
Classez les indicateurs suivants selon le niveau CTI et le type
(IoC / IoA / TTP)~: (a)~\texttt{sha256:3a7d4f8...}~;
(b)~exécution de \texttt{powershell.exe -enc ...} à 03h00~;
(c)~utilisation de l'outil Mimikatz~; (d)~campagne de phishing
ciblant le secteur bancaire africain~; (e)~domaine \texttt{evil.io}.

\medskip\textbf{Ex.~9.2 — Objet STIX.}
Rédigez en JSON STIX~2.1 un objet \texttt{indicator} pour le domaine
\texttt{evil.io} utilisé comme C2 lors de l'incident LiZbiZ-SIS.
Complétez avec un objet \texttt{relationship} liant cet indicateur
à la technique ATT\&CK T1071.001.

\medskip\textbf{Ex.~9.3 — Intégration SIEM.}
Décrivez le processus complet d'intégration d'un flux MISP dans
Wazuh pour LiZbiZ-SIS~: script d'export, format CDB, règle Wazuh,
alerte générée, fréquence de mise à jour recommandée.
\end{boxexo}
